\documentclass[12pt,a4paper]{article}

\usepackage[utf8]{inputenc}
\usepackage[brazil]{babel}
\usepackage[T1]{fontenc}
\usepackage{geometry}
\usepackage{setspace}
\usepackage{indentfirst}
\usepackage{listings}
\usepackage{xcolor}

\geometry{margin=2.5cm}
\onehalfspacing

\lstset{
    basicstyle=\ttfamily\small,
    breaklines=true,
    frame=single,
    numbers=left,
    numberstyle=\tiny,
    keywordstyle=\bfseries,
    showstringspaces=false
}

\begin{document}

\begin{titlepage}
\begin{center}

\textbf{CENTRO UNIVERSITÁRIO MAURÍCIO DE NASSAU}

\vspace{2.5cm}

\textbf{\Large DESENVOLVIMENTO DO COMPILADOR LUDICODE}

\vspace{2cm}

\textbf{Equipe:}

\vspace{0.5cm}

Abraão Araujo dos Santos\\
Erick Alves de Souza\\
Evandro França da Silva Filho\\
Gabriel Marques Barbosa de Santana\\
Gabriella Maria Nascimento da Silva\\
Jonnas Kauan Pereira Novaes\\
Mateus De Miranda Santos Moura\\
Matheus Pinheiro da Cunha\\

\vfill

Maceió - AL\\
2026

\end{center}
\end{titlepage}

\tableofcontents
\newpage

\section{Introdução}

O projeto LudiCode foi desenvolvido com o objetivo de criar uma linguagem de programação voltada para o ensino de lógica computacional e robótica. A proposta da linguagem é utilizar comandos em português, facilitando o entendimento por parte de estudantes que estão tendo os primeiros contatos com programação.

A implementação do projeto foi feita utilizando a linguagem Python e a biblioteca SLY, que auxilia na construção de analisadores léxicos e sintáticos. O sistema segue uma estrutura semelhante à utilizada em compiladores e interpretadores, passando por etapas como análise léxica, análise sintática, geração da árvore sintática abstrata e execução dos comandos.

Dessa forma, o projeto permite relacionar os conceitos estudados na disciplina de Compiladores com uma aplicação prática, demonstrando como uma linguagem de programação pode ser definida, analisada e executada.

\section{Estrutura da Linguagem}

A linguagem LudiCode possui uma sintaxe simples e próxima da linguagem natural. Ela permite o uso de variáveis, estruturas condicionais, laços de repetição, operações matemáticas e comandos de exibição.

Um exemplo de código escrito na linguagem é apresentado a seguir:

\begin{lstlisting}
var idade = 18

se idade >= 18 {
    mostrar("Maior de idade")
}
\end{lstlisting}

Nesse exemplo, uma variável chamada \texttt{idade} é criada com o valor 18. Em seguida, é feita uma verificação condicional para exibir uma mensagem caso a condição seja verdadeira.

A utilização de palavras em português torna a linguagem mais acessível, principalmente para estudantes iniciantes, pois reduz a dificuldade inicial causada por comandos em inglês.

\section{Gramática da Linguagem}

A gramática da linguagem define quais comandos são considerados válidos. No LudiCode, essas regras estão relacionadas à forma como variáveis, expressões, condições e blocos de comandos devem ser escritos.

De forma simplificada, uma declaração de variável pode ser representada da seguinte maneira:

\begin{lstlisting}
var identificador = expressao
\end{lstlisting}

Uma estrutura condicional segue o formato:

\begin{lstlisting}
se condicao {
    comandos
}
\end{lstlisting}

E uma estrutura de repetição pode ser representada como:

\begin{lstlisting}
enquanto condicao {
    comandos
}
\end{lstlisting}

Essas regras são importantes porque permitem que o analisador sintático reconheça se o programa foi escrito corretamente. Caso o código não siga a estrutura esperada, o sistema pode identificar e informar o erro.

\section{Desenvolvimento do Sistema}

O sistema foi organizado em módulos, cada um responsável por uma etapa do processamento do código-fonte. Essa divisão facilita a compreensão do projeto e torna o código mais organizado.

O fluxo geral do sistema pode ser representado da seguinte forma:

\begin{center}
Código Fonte $\rightarrow$ Lexer $\rightarrow$ Parser $\rightarrow$ AST $\rightarrow$ Interpretador
\end{center}

Primeiro, o código escrito pelo usuário é lido pelo analisador léxico. Depois, os tokens gerados são enviados para o analisador sintático. Em seguida, é construída uma Árvore Sintática Abstrata, que será utilizada pelo interpretador para executar o programa.

\subsection{Arquivo main.py}

O arquivo \texttt{main.py} é responsável por iniciar a execução do sistema. Ele realiza a leitura do código-fonte e aciona os demais componentes do compilador.

De maneira geral, esse arquivo funciona como o ponto de entrada do projeto. Ele recebe o programa escrito na linguagem LudiCode, envia esse conteúdo para o lexer, depois para o parser e, por fim, executa o interpretador.

Essa organização permite que as etapas do compilador sejam executadas em sequência, garantindo que o código passe por validações antes da execução.

\subsection{Analisador Léxico}

O analisador léxico foi implementado no arquivo \texttt{lexer.py}. Sua função é ler o código-fonte e identificar os elementos básicos da linguagem, chamados de tokens.

Tokens são partes reconhecidas do código, como palavras reservadas, identificadores, números, operadores e símbolos. Por exemplo, na instrução:

\begin{lstlisting}
var idade = 18
\end{lstlisting}

o analisador identifica elementos como a palavra reservada \texttt{var}, o identificador \texttt{idade}, o operador de atribuição e o número \texttt{18}.

Entre as principais palavras reservadas reconhecidas pela linguagem estão:

\begin{itemize}
    \item \texttt{var};
    \item \texttt{se};
    \item \texttt{senao};
    \item \texttt{enquanto};
    \item \texttt{mostrar};
    \item \texttt{verdadeiro};
    \item \texttt{falso}.
\end{itemize}

Essa etapa é fundamental, pois transforma o texto digitado pelo usuário em uma sequência de informações que pode ser processada pelas próximas fases.

\subsection{Analisador Sintático}

A análise sintática é realizada pelo arquivo \texttt{parse.py}. Essa etapa verifica se a sequência de tokens gerada pelo analisador léxico está de acordo com as regras da gramática da linguagem.

Por exemplo, a instrução abaixo segue a estrutura correta da linguagem:

\begin{lstlisting}
var nome = "LudiCode"
\end{lstlisting}

Por outro lado, uma instrução com a ordem incorreta dos elementos pode gerar um erro sintático:

\begin{lstlisting}
var = nome "LudiCode"
\end{lstlisting}

Além de verificar a sintaxe, o parser também é responsável por construir a Árvore Sintática Abstrata, conhecida como AST.

\subsection{Árvore Sintática Abstrata}

A Árvore Sintática Abstrata, ou AST, é uma estrutura utilizada para representar o programa de forma organizada. Em vez de executar diretamente o texto escrito pelo usuário, o sistema cria uma representação interna mais fácil de processar.

Por exemplo, uma declaração de variável pode ser representada internamente como uma estrutura contendo o tipo da instrução, o nome da variável e o valor atribuído.

Essa representação facilita a execução do programa, pois o interpretador consegue percorrer a árvore e executar cada comando na ordem correta.

\subsection{Interpretador}

O interpretador foi implementado no arquivo \texttt{interpreter.py}. Ele é responsável por percorrer a AST e executar as instruções representadas nela.

Durante a execução, o interpretador realiza operações como:

\begin{itemize}
    \item criação e atualização de variáveis;
    \item avaliação de expressões matemáticas;
    \item verificação de condições;
    \item execução de laços de repetição;
    \item exibição de mensagens;
    \item tratamento de erros.
\end{itemize}

O interpretador também utiliza uma tabela de símbolos, que funciona como uma área de armazenamento para as variáveis criadas durante a execução do programa. Assim, quando uma variável é declarada, seu nome e valor ficam guardados para uso posterior.

\section{Tratamento de Erros}

O tratamento de erros é uma parte importante do projeto, pois ajuda o usuário a identificar problemas no código.

Entre os erros que podem ser tratados estão:

\begin{itemize}
    \item erros de sintaxe;
    \item uso de variáveis não declaradas;
    \item divisão por zero;
    \item operações inválidas entre tipos diferentes;
    \item comandos escritos de forma incorreta.
\end{itemize}

Essas verificações tornam o sistema mais seguro e facilitam o aprendizado, pois o usuário consegue entender melhor onde está o problema no código desenvolvido.

\section{Testes Realizados}

Para verificar o funcionamento da linguagem, foram utilizados arquivos de teste presentes no projeto. Esses testes avaliam diferentes partes do sistema, garantindo que os módulos trabalhem de forma integrada.

Entre os principais cenários avaliados estão:

\begin{itemize}
    \item declaração de variáveis;
    \item operações matemáticas;
    \item estruturas condicionais;
    \item estruturas de repetição;
    \item validação de tipos;
    \item divisão por zero.
\end{itemize}

Os testes demonstram que o sistema consegue processar diferentes comandos da linguagem e identificar erros durante a execução. Isso mostra que as etapas de análise léxica, análise sintática e interpretação estão funcionando em conjunto.

\section{Conclusão}

O desenvolvimento do LudiCode permitiu aplicar na prática conceitos estudados na disciplina de Compiladores. O projeto mostra como uma linguagem de programação pode ser criada a partir da definição de tokens, regras gramaticais e mecanismos de execução.

Durante o trabalho, foram implementadas etapas fundamentais, como análise léxica, análise sintática, construção da AST e interpretação dos comandos. Cada uma dessas etapas possui uma função específica dentro do funcionamento do sistema.

Além do aspecto técnico, o projeto também possui uma proposta educacional, pois utiliza comandos em português para facilitar o aprendizado de programação. Dessa forma, o LudiCode demonstra como conceitos de compiladores podem ser aplicados em uma ferramenta voltada ao ensino de lógica computacional e robótica.

\end{document}
